\section{Related Work}\label{sec:related}

For each line of work we state what it does, what it shares with this study and how this study differs.

\subsection{Optimistic replication and offline work}
Saito and Shapiro survey \emph{optimistic replication} and separate the detection of concurrent updates from their
resolution~\cite{saito2005optimistic}, the split we use between the clocks of a strategy and its merge rule. The Coda file system made
\emph{disconnected operation} practical for mobile clients: a laptop hoards files, works on them offline and
reintegrates the changes, flagging write/write conflicts for manual repair~\cite{kistler1992coda}. Coda works at the
granularity of whole files, which is the record-level granularity of our LWW baseline. Bayou lets applications attach a
dependency check and a merge procedure to every write, with an order made final at a primary
server~\cite{terry1995bayou}; our strategies need no such commit order. Dynamo uses vector clocks to keep
concurrent versions of a key as siblings and returns them to the application for semantic
reconciliation~\cite{decandia2007dynamo}. Our version-vector strategy follows the same
idea with a field-level merge and the mutual-inconsistency detection of Parker et al.~\cite{parker1983detection}.
These systems describe mechanisms; none of them measures how many user edits each mechanism loses on one workload.

\subsection{CRDT implementations and their cost}
CRDTs remove the need for application merge code by restricting the data types to ones whose concurrent updates
commute~\cite{shapiro2011crdt,shapiro2011comprehensive,preguica2018overview}. The general-purpose designs pay for
this with metadata. Attiya et al.\ show that a replicated list with the expected behaviour needs metadata that grows
with the number of operations~\cite{attiya2016editing}, which explains the history kept by Automerge and
Yjs~\cite{kleppmann2017json,nicolaescu2016yjs}. Two lines of work reduce the cost: delta-state
CRDTs~\cite{almeida2018delta} ship only recent changes, and optimised observed-remove sets drop the tombstones of
removed elements by keeping a compact causal context~\cite{bieniusa2012orset}. Convergence proofs have been
mechanised~\cite{gomes2017verifying}, and the \emph{crdt-benchmarks} suite compares the encoding
size and speed of CRDT libraries on synthetic and real text-editing traces~\cite{crdtbenchmarks}. Our hybrid strategy
reuses the delta-state model and the causal contexts of these works. The difference is the workload: published
benchmarks and proofs centre on text and lists, while a typical mobile application stores records of scalar fields
with a few small collections, where most of the history of a sequence CRDT is never needed.

\subsection{Synchronisation frameworks for mobile applications}
Local-first software~\cite{kleppmann2019localfirst} argues for applications that keep the primary copy of the data
on the device and synchronise it with CRDTs; its authors built prototypes on Automerge and name the performance and
storage cost of the history among the open problems. Automerge~\cite{automergedocs} and Yjs~\cite{yjsdocs} are the two
widely used libraries: Automerge offers a JSON document with a sync protocol based on Bloom filters of change hashes,
Yjs offers shared maps, arrays and text with a state-vector exchange. Automerge's core is compiled to WebAssembly,
which the default Hermes engine of React Native does not run; we met exactly this limitation
(Section~\ref{sec:setup}). WatermelonDB~\cite{watermelondbdocs}, a database for React Native, synchronises through a
pull/push protocol with per-record change tracking; when a record changed on both sides, the locally changed columns
win (``per-column client-wins''). Such frameworks, like the textbooks~\cite{kleppmann2017ddia,vansteen2017distributed},
make or discuss the policy choice qualitatively; our study measures what it costs.

\subsection{Gap analysis}
The works above provide the building blocks and their guarantees, but no study compares LWW, version vectors and CRDT libraries on the same mobile application, with
structured records instead of text, with clock skew, and with a precise count of lost edits. We address this gap with
a common strategy interface, a ground-truth oracle and a type-aware hybrid that applies the cheap building blocks per
field. Table~\ref{tab:expectations} states the expectations for the hybrid that the experiments must confirm or
refute. To validate the measurement harness we use public data: the workloads of the crdt-benchmarks suite,
including its real editing trace, will be replayed through our Automerge and Yjs adapters and the byte counts compared
with the published ones (Section~\ref{sec:validation}).

\begin{table}[ht]
\caption{Expected position of the hybrid strategy (S4) relative to the baselines, per metric.}
\label{tab:expectations}
\small
\begin{tabular}{@{}lll@{}}
\toprule
Metric & Expectation for S4 & Experiments \\
\midrule
Convergence & always, as for CRDTs & all \\
Lost operations & far below LWW, close to a full CRDT & E1, E2, E4, E5 \\
Payload and state size & below a full CRDT, near LWW & E3, E5 \\
Implementation effort & between LWW and a CRDT library & E6 \\
\bottomrule
\end{tabular}
\end{table}
